\item Two touching bars 1 and 2 are placed on an inclined plane forming an angle \(\alpha\) with the horizontal (Fig. 1.10). The masses of the bars are equal to \(m_1\) and \(m_2\), and the coefficients of friction between the inclined plane and these bars are equal to \(k_1\) and \(k_2\) respectively, with \(k_1 > k_2\). Find:
    \begin{center}
        \begin{tikzpicture}
            \node at (0, 0) {{image1.png}};
        \end{tikzpicture}
    \end{center}
    \begin{enumerate}
        \item the force of interaction of the bars in the process of motion;
        \item the minimum value of the angle \(\alpha\) at which the bars start sliding down.
    \end{enumerate}

    \begin{solution}
        \begin{center}
            \begin{tikzpicture}
                \pic at (0, 0) {frame=3cm};
            \end{tikzpicture}
        \end{center}
        
        \begin{enumerate}
            \item 
            \begin{align*}
                \intertext{Let us indicate the positive direction of \( x \)-axis along the incline (see figure). The figures show the force diagram for the blocks.}
                \intertext{Let \( R \) be the force of interaction between the bars and they are obviously sliding down with the same constant acceleration \( w \).}
                \intertext{Newton's second law of motion in projection form along \( x \)-axis for the blocks gives}
                m_1 g \sin \alpha - k_1 m_1 g \cos \alpha + R &= m_1 w \tag{1}\\
                m_2 g \sin \alpha - R - k_2 m_2 g \cos \alpha &= m_2 w \tag{2}
                \intertext{Solving Eqs. (1) and (2) simultaneously, we get}
                w &= g \sin \alpha - g \cos \alpha \dfrac{k_1 m_1 + k_2 m_2}{m_1 + m_2} \tag{3}\\
                \intertext{and}
                R &= \dfrac{m_1 m_2 (k_1 - k_2) g \cos \alpha}{m_1 + m_2}            
            \end{align*}
        \item 
            \begin{align*}
                \intertext{When the blocks just slide down the plane, \( w = 0 \), then from Eq. (3)}
                g \sin \alpha - g \cos \alpha \dfrac{k_1 m_1 + k_2 m_2}{m_1 + m_2} &= 0\\
                \intertext{or}
                (m_1 + m_2) \sin \alpha &= (k_1 m_1 + k_2 m_2) \cos \alpha
                \intertext{Hence,}
                \tan \alpha &= \dfrac{k_1 m_1 + k_2 m_2}{m_1 + m_2}
            \end{align*}
    \end{enumerate}
    \end{solution}